% Einheit 2 von 3 – Zehnerpotenzen (bank/potenzen-wurzeln/e2.jsonl, zusammenbau v0.1)
%% TODO Zweigzeile Teil 1: was hier gelernt wird (Satz in Schülersprache aus dem Einheitstitel) – Einheit 2
\einheitenkopf[e2]{Einheit 2 von 3 · Zehnerpotenzen}
\zweigzeile{ab Kl. 9, je nach Buch bis Kl. 10 (am Gymnasium ab Kl. 8) · P10}
\setcounter{aufgabe}{18}

%% TODO Titel: Ich-kann-Satz fehlt in der Bank (Platzhalter: Kettenname) – Nr. 19
%% TODO Anweisung: ein Satz über der ersten Teilaufgabe fehlt in der Bank – Nr. 19
\begin{aufgabe}{Zehnerpotenzen}
%% TODO \rechenplatz{n}: Schrittzahl des Grundfalls fehlt in der Bank (nur bei fester Schreibform, 2.3 b)
\begin{teile}
\teil $4{,}7 \cdot 10^{8}$ – größer oder kleiner als eins?\\ \kreuz{größer als eins} \\ \kreuz{kleiner als eins}
\teil $10^{4}$ – ausgeschrieben? \leerfeld
\teil $10^{9}$ – welches Zahlwort?\\ \kreuz{Million} \\ \kreuz{Milliarde} \\ \kreuz{Billion}
\teil $53 \cdot 100$ – wie viel? \leerfeld
\teil $29 \cdot 10^{5}$ – wie viel? \leerfeld
\teil $4{,}8 \cdot 10^{3}$ – wie viel? \leerfeld
\teil $6{,}3 \cdot 10^{5}$ – ausgeschrieben? \hfill (P10 2016 OS) \\ \leerfeld
\teil $58\,000$ – in Zehnerpotenzschreibweise? \leerfeld
\teil $64\,000 = 6{,}4 \cdot 10^{\square}$ – welche Hochzahl? \leerfeld
\teil Das Gehirn eines Menschen hat etwa $8{,}6 \cdot 10^{10}$ Nervenzellen. Welche Zahl passt? Kreuze an. \hfill (P10 2015 OS)\\ \kreuz{$860\,000\,000\,000$} \\ \kreuz{$86\,000\,000\,000$} \\ \kreuz{$8\,600\,000\,000$}
\teil $3{,}6 \cdot 10^{-3}$ – als Dezimalzahl? \leerfeld
\teil $0{,}0052$ – in Zehnerpotenzschreibweise? \leerfeld
\teil $4{,}1 \cdot 10^{5}$; $9{,}3 \cdot 10^{4}$; $2{,}8 \cdot 10^{6}$ – aufsteigend geordnet?
\teil Der Taschenrechner zeigt 6.2E7 – welche Zahl ist das? \leerfeld
\teil Von der Erde zur Sonne sind es etwa $149\,597\,871$ km – in Zehnerpotenzschreibweise, Faktor auf zwei Stellen nach dem Komma gerundet? \leerfeld
\teil Ein Bakterium ist etwa $2 \cdot 10^{-3}$ mm lang. Wie viele liegen hintereinander auf $1$ mm? \leerfeld Bakterien
\teil $10^{4} \cdot 10^{3}$ – als Zehnerpotenz? \leerfeld
\teil $730\,000 = 7{,}3 \cdot 10^{\square}$ – welche Hochzahl? \hfill (P10 2025 OS) \\ \leerfeld
\end{teile}
\end{aufgabe}

%% TODO Titel: Ich-kann-Satz fehlt in der Bank (Platzhalter: Kettenname) – Nr. 20
%% TODO Anweisung: ein Satz über der ersten Teilaufgabe fehlt in der Bank – Nr. 20
\begin{aufgabe}{Zehnerpotenzen – Fehler finden, Begründen}
\begin{teile}
\teil Ben schreibt $920\,000 = 9{,}2 \cdot 10^{4}$. Wo ist der Fehler?
\teil Begründe, warum $2{,}6 \cdot 10^{5}$ und $26 \cdot 10^{4}$ dieselbe Zahl sind.
\end{teile}
\end{aufgabe}
